\documentclass[12pt,letterpaper,oneside]{report}

% Kompilator w Overleaf: pdfLaTeX
\usepackage[utf8]{inputenc}
\usepackage[T1]{fontenc}
\usepackage[polish]{babel}
\usepackage{polski}

\usepackage[letterpaper,left=2.0cm,right=2.0cm,top=2.0cm,bottom=2.0cm]{geometry}
\usepackage{microtype}
\setlength{\parindent}{1.25em}
\setlength{\parskip}{0pt}
\setlength{\emergencystretch}{2em}

\usepackage{amsmath,amssymb,mathtools,bm}
\usepackage{graphicx}
\graphicspath{{figures/}}
\usepackage{booktabs,longtable,array,calc}
\usepackage{etoolbox}
\usepackage{caption}
\captionsetup{font=small,labelfont=normalfont,justification=justified,singlelinecheck=false}
\usepackage{float}
\usepackage{enumitem}
\setlist[itemize]{topsep=0.35em,itemsep=0.15em,leftmargin=1.6em}
\setlist[enumerate]{topsep=0.35em,itemsep=0.2em,leftmargin=1.9em}

\usepackage{titlesec}
\titleformat{\chapter}[display]
  {\normalfont\Large\bfseries}
  {\chaptername\ \thechapter}
  {0.85em}
  {}
\titlespacing*{\chapter}{0pt}{0pt}{2.2em}
\titleformat{\section}
  {\normalfont\large\bfseries}
  {\thesection}
  {0.9em}
  {}
\titlespacing*{\section}{0pt}{1.7em}{0.65em}
\titleformat{\subsection}
  {\normalfont\normalsize\bfseries}
  {\thesubsection}
  {0.8em}
  {}
\titlespacing*{\subsection}{0pt}{1.3em}{0.45em}

% W pracy wzorcowej numeracja podrozdziałów zaczyna się od 1 w każdym rozdziale.
\renewcommand{\thesection}{\arabic{section}}
\renewcommand{\thesubsection}{\thesection.\arabic{subsection}}
\numberwithin{equation}{chapter}
\numberwithin{figure}{chapter}
\numberwithin{table}{chapter}

\usepackage{tocloft}
\renewcommand{\contentsname}{Spis treści}
\renewcommand{\cfttoctitlefont}{\hfill\Large\bfseries}
\renewcommand{\cftaftertoctitle}{\hfill}
\setlength{\cftbeforetoctitleskip}{0.5cm}
\setlength{\cftaftertoctitleskip}{1.0cm}
\setlength{\cftbeforechapskip}{0.25em}
\setlength{\cftbeforesecskip}{0pt}

% Dane strony tytułowej - wystarczy zmienić poniższe linie.
\newcommand{\UniversityName}{Uniwersytet Jagielloński w Krakowie}
\newcommand{\FacultyName}{Wydział Fizyki, Astronomii i Informatyki Stosowanej}
\newcommand{\AuthorName}{Karina Nepravskaya}
\newcommand{\AlbumNumber}{1209730}
\newcommand{\SupervisorName}{tytuł, imię i nazwisko promotora}
\newcommand{\DepartmentName}{nazwa zakładu lub instytutu}
\newcommand{\YearOfSubmission}{2026}

\usepackage{hyperref}
\hypersetup{
  colorlinks=true,
  linkcolor=blue,
  citecolor=blue,
  urlcolor=blue,
  pdftitle={Analiza korelacji pędów elektronów w potrójnej jonizacji neonu},
  pdfauthor={\AuthorName}
}

% Skróty używane w tekście.
\newcommand{\Pe}{\ensuremath{P_e}}
\newcommand{\Neff}{\ensuremath{N_{\mathrm{eff}}}}
\newcommand{\xmax}{\ensuremath{x_{\max}}}
\newcommand{\direct}{\textit{direct}}
\newcommand{\delayed}{\textit{delayed}}

% Zwykły, pozbawiony ramki akapit wyróżniony - zgodny z oszczędnym stylem pracy wzorcowej.
\newenvironment{importantbox}[1]
  {\par\medskip\noindent\textbf{#1.}\quad}
  {\par\medskip}

% Pomocnicze definicje wymagane przez tabele wygenerowane przez Pandoc.
\providecommand{\tightlist}{\setlength{\itemsep}{0pt}\setlength{\parskip}{0pt}}
\makeatletter
\patchcmd\longtable{\par}{\if@noskipsec\mbox{}\fi\par}{}{}
\makeatother

\pagestyle{plain}
\setcounter{secnumdepth}{2}
\setcounter{tocdepth}{2}

\begin{document}

\begin{titlepage}
\thispagestyle{empty}

\noindent
{\bfseries \UniversityName}\par
\noindent
\FacultyName

\vspace{2.7cm}

\begin{center}
\centering
\LARGE Karina Nepravskaya

\vspace{0.35cm}
{\large Nr albumu: \AlbumNumber}
\end{center}

\vspace{2.0cm}

\begin{center}
{\LARGE\bfseries
Analiza korelacji pędów elektronów\\[0.2em]
w potrójnej jonizacji neonu\\[0.2em]
w silnym polu laserowym\par}

\vspace{1.2cm}

{\normalsize Notatka metodologiczna}
\end{center}

\vfill

\begin{flushright}
\hfill
\begin{minipage}{0.48\textwidth}
\raggedleft
Praca wykonana pod kierunkiem\\
\SupervisorName\\
\DepartmentName
\end{minipage}
\end{flushright}

\vfill

\begin{center}
Kraków \YearOfSubmission
\end{center}
\end{titlepage}

\pagenumbering{arabic}
\setcounter{page}{2}
\tableofcontents
\clearpage
\phantomsection
\addcontentsline{toc}{chapter}{Streszczenie}
\begin{center}
  {\bfseries Streszczenie}
\end{center}
\vspace{1.2em}

Celem analizy jest wyodrębnienie i ilościowe opisanie korelacji między
końcowymi składowymi pędu trzech elektronów emitowanych w potrójnej
jonizacji atomu neonu. Punktem wyjścia była analiza głównych składowych
przeprowadzona w dziewięciowymiarowej przestrzeni składowych pędu. PCA
wskazała, że dynamikę wzdłuż osi polaryzacji lasera można rozdzielić na
jeden kierunek kolektywny, odpowiadający zgodnemu ruchowi wszystkich
elektronów, oraz dwa niezależne kierunki względne, opisujące
nierównomierność podziału pędu. Zamiast posługiwać się bezpośrednio
zależnymi od zbioru danymi wektorami własnymi PCA, wprowadzono stałą,
ortonormalną bazę Jacobiego--Helmerta. Umożliwia ona porównywanie
kanałów direct i delayed w tym samym układzie współrzędnych oraz
zachowuje pełną informację o trzech podłużnych składowych pędu.

Współrzędne liniowe uzupełniono obserwablami zaczerpniętymi z
reprezentacji simpleksowej: udziałami bezwzględnego pędu, efektywną
liczbą elektronów uczestniczących w podziale pędu, największym udziałem
jednego elektronu, entropią podziału, sektorami znaków oraz
współrzędnymi logarytmicznych ilorazów. Taki zestaw oddziela trzy różne
aspekty zdarzenia: skalę i kierunek ruchu kolektywnego, stopień
nierównowagi względnej oraz strukturę podziału wartości bezwzględnych i
znaków pędów. Dla danych przy natężeniu \ensuremath{1{,}7\,\mathrm{PW/cm^2}} kanał direct
charakteryzuje się silniejszymi dodatnimi korelacjami parami, mniejszą
wartością względnej nierównowagi K i bardziej równomiernym podziałem
pędu niż kanał delayed. Wyniki te są zgodne z obrazem bardziej
kolektywnej, niemal jednoczesnej emisji w kanale direct oraz bardziej
rozdzielonej czasowo dynamiki w kanale delayed, przy czym same pędy
końcowe nie stanowią samodzielnego dowodu mechanizmu czasowego.

\begin{importantbox}{Najważniejsza idea metody}
PCA służy tutaj jako narzędzie rozpoznania struktury danych, a nie jako
ostateczny układ współrzędnych. Ostateczna baza jest ustalona przez
kinematykę układu wielu cząstek. Dzięki temu osie nie obracają się przy
zmianie kanału jonizacji, natężenia pola ani liczebności próbki.
\end{importantbox}


\clearpage
\thispagestyle{empty}
\null
\clearpage
\chapter{Cel pracy i dane}

\section{Pytanie badawcze}

W potrójnej jonizacji neonu interesuje nas nie tylko rozkład pędu
pojedynczego elektronu, lecz także sposób, w jaki pęd końcowy dzieli się
między trzy elektrony w jednym zdarzeniu. Celem pracy jest zbudowanie
opisu, który rozdziela ruch wspólny elektronów od ich wzajemnej
nierównowagi, a następnie porównanie kanałów \textit{direct} i
\textit{delayed}.

Analiza dotyczy składowych pędu wzdłuż osi polaryzacji lasera. W tym
kierunku pole nadaje elektronom dominujący pęd dryfu, dlatego dalsze
rozważania prowadzone są dla
\[
p_2\equiv p_{2z},\qquad p_3\equiv p_{3z},\qquad p_4\equiv p_{4z}.
\]
Elektrony 2 i 3 są w danych trajektoryjnych oznaczone jako początkowo
związane, a elektron 4 jako początkowo tunelujący. Są to etykiety
historii trajektorii; obserwable zależne od numeru elektronu trzeba więc
odróżniać od wielkości niezmienniczych na permutacje.

\section{Dlaczego nie wystarcza simpleks}

Dla trzech elektronów udziały bezwzględnego pędu
\begin{equation}
x_i=\frac{|p_i|}{|p_2|+|p_3|+|p_4|},\qquad x_2+x_3+x_4=1,
\label{eq:shares-main1}
\end{equation}
można przedstawić w trójkącie simpleksowym. Reprezentacja ta dobrze
opisuje \emph{względny} podział wartości bezwzględnych, ale nie zawiera
znaków pędów ani ich bezwzględnej skali. Ponadto dla większej liczby
elektronów simpleks staje się obiektem wyższowymiarowym i przestaje być
czytelną pojedynczą wizualizacją.

W pracy zastępujemy więc jeden przeciążony wykres zestawem obserwabli o
oddzielnych rolach: współrzędnymi ruchu kolektywnego i względnego,
miarami równomierności udziałów oraz sektorami znaków. Dzięki temu każda
liczba i każdy wykres odpowiadają na jedno konkretne pytanie.

\chapter{Rozdzielenie ruchu kolektywnego i względnego}

\section{Rola PCA}

Najpierw wykonano PCA w dziewięciowymiarowej przestrzeni końcowych
składowych pędu elektronów. Dane standaryzowano, dlatego PCA bada
strukturę korelacji, a nie wyłącznie różnice wariancji między kolumnami.
Dla kanału \textit{direct} pierwsza składowa wyjaśnia około \(30{,}0\%\)
całkowitej wariancji i ma niemal jednakowe współczynniki przy
\(p_{2z},p_{3z},p_{4z}\):
\[
\mathrm{PC1}\approx
0{,}580p_{2z}+0{,}577p_{3z}+0{,}572p_{4z}.
\]
Wskazuje to na kierunek zbliżony do \((1,1,1)\), czyli na wspólny ruch
trzech elektronów. Dwa pozostałe niezależne kierunki podłużne opisują
różnice pomiędzy elektronami.

PCA jest tu narzędziem diagnostycznym, nie ostatecznym układem
współrzędnych. Wektory własne zależą od konkretnej próbki; ich znaki są
umowne, a osie o podobnych wartościach własnych mogą obracać się w tej
samej podprzestrzeni. Bezpośrednie porównywanie osi PCA między kanałami
byłoby zatem niejednoznaczne.

\section{Stała baza Jacobiego--Helmerta}

Dla wektora \(\mathbf p_z=(p_2,p_3,p_4)^{\mathsf T}\) wprowadzamy
ortonormalną bazę
\begin{equation}
Q=\frac{p_2+p_3+p_4}{\sqrt3},\qquad
\xi_1=\frac{p_2-p_3}{\sqrt2},\qquad
\xi_2=\frac{p_2+p_3-2p_4}{\sqrt6}.
\label{eq:helmert-main1}
\end{equation}
Współrzędna \(Q\) opisuje wspólny podłużny ruch elektronów.
Współrzędne \(\xi_1\) i \(\xi_2\) opisują dwa niezależne kontrasty
wewnętrzne: różnicę w parze (2,3) oraz ruch tej pary względem elektronu
4. Układ jest stały dla wszystkich analizowanych zbiorów, dlatego
porównanie \textit{direct}--\textit{delayed} nie zależy od przypadkowej
orientacji osi PCA.

Transformacja jest zmianą bazy, a nie redukcją danych:
\begin{equation}
p_2^2+p_3^2+p_4^2=Q^2+\xi_1^2+\xi_2^2.
\label{eq:norm-main1}
\end{equation}
Zachowuje więc pełną informację o trzech podłużnych pędach. Do
wizualizacji ruchu kolektywnego stosujemy także
\(\Pe{}=p_2+p_3+p_4=\sqrt3\,Q\).

\section{Miary nierównowagi i podziału}

Promień w płaszczyźnie względnej
\begin{equation}
K=\sqrt{\xi_1^2+\xi_2^2}
\label{eq:k-main1}
\end{equation}
mierzy wielkość nierównowagi pędów. Jest równy zeru wtedy i tylko wtedy,
gdy \(p_2=p_3=p_4\), oraz nie zależy od permutacji elektronów:
\[
K^2=\sum_{i=2}^{4}(p_i-\bar p)^2,\qquad
\bar p=\frac{p_2+p_3+p_4}{3}.
\]
Samo \(K\) nie mówi jednak, który kontrast dominuje; tę informację
zachowuje para \((\xi_1,\xi_2)\).

Udziały z równania~\eqref{eq:shares-main1} opisujemy trzema
uzupełniającymi miarami:
\[
N_{\mathrm{eff}}=\frac{1}{x_2^2+x_3^2+x_4^2},\qquad
x_{\max}=\max_i x_i,\qquad
H=-\frac{\sum_i x_i\ln x_i}{\ln3}.
\]
Wartości \(N_{\mathrm{eff}}\) bliskie 3, małe \(x_{\max}\) i \(H\)
bliskie 1 oznaczają równomierny podział wartości bezwzględnych.
Informację o kierunku emisji przywracają osiem sektorów znaków
\((\operatorname{sgn}p_2,\operatorname{sgn}p_3,\operatorname{sgn}p_4)\).

\chapter{Wyniki dla \(I=1{,}7\,\mathrm{PW/cm^2}\)}

\section{Porównanie kanałów}

Analiza obejmuje 1695 zdarzeń \textit{direct} i 1601 zdarzeń
\textit{delayed}. Najważniejsze wielkości zebrano w tabeli
\ref{tab:results-main1}. Podajemy średnie i odchylenia standardowe;
porównując kształty histogramów, należy używać tych samych granic binów
oraz normalizacji do gęstości prawdopodobieństwa.

\begin{table}[htbp]
\centering
\caption{Zwięzłe porównanie kanałów jonizacji.}
\label{tab:results-main1}
\begin{tabular}{lrrrrrr}
\toprule
kanał & \(n\) & \(\sigma(\Pe{})\) & \(\langle K\rangle\) &
\(\langle N_{\mathrm{eff}}\rangle\) & \(\langle x_{\max}\rangle\) &
\(r_{23}\)\\
\midrule
direct  & 1695 & 7,543 & 1,398 & 2,628 & 0,473 & 0,868\\
delayed & 1601 & 5,384 & 2,383 & 2,374 & 0,533 & 0,214\\
\bottomrule
\end{tabular}
\end{table}

Kanał \textit{direct} ma mniejsze \(\langle K\rangle\), większe
\(N_{\mathrm{eff}}\) i mniejsze \(x_{\max}\). Te trzy niezależne miary
mówią to samo: wartości bezwzględne pędu są w nim dzielone bardziej
równomiernie. Wysoka korelacja \(r_{23}=0{,}868\) jest spójna z
kolektywnym ruchem elektronów. Dla kanału \textit{delayed} rozkład
względny jest szerszy, korelacje są słabsze, a pojedynczy elektron
częściej przejmuje większą część pędu.

Sektory znaków wzmacniają ten wniosek. W kanale \textit{direct} sektor
\(+++\) obejmuje \(47{,}4\%\) zdarzeń, a sektory zgodnego kierunku
\((+++)\) i \((---)\) dominują nad sektorami mieszanymi. W kanale
\textit{delayed} udział sektora \(+++\) spada do \(23{,}0\%\), a emisja
o mieszanych znakach jest wyraźnie częstsza. Nie jest to informacja
zawarta w samych udziałach \(|p_i|\), dlatego sektory znaków pozostają
koniecznym uzupełnieniem analizy simpleksowej.

\section{Interpretacja}

Wyniki są zgodne z obrazem bardziej zbiorowej emisji w kanale
\textit{direct}: elektrony opuszczają obszar jonu w zbliżonej fazie pola,
co prowadzi do podobnych znaków i silnie dodatnich korelacji pędów.
Kanał \textit{delayed} jest zgodny z emisją rozdzieloną czasowo, w której
co najmniej jeden elektron doświadcza przyspieszania w innej fazie pola.

To nie jest samodzielny dowód mechanizmu czasowego. Końcowe pędy są
projekcją całej dynamiki na stan asymptotyczny. Wniosek ma więc charakter
ilościowej zgodności rozkładów pędu z klasyfikacją trajektorii, a nie
zastępuje analizy czasów jonizacji.

\chapter{Zakres stosowalności i wnioski}

Zaproponowany opis ma trzy zalety. Po pierwsze, transformacja
Jacobiego--Helmerta zachowuje pełną informację o trzech podłużnych
pędach. Po drugie, jej osie są ustalone przez kinematykę, a nie przez
konkretną próbkę, dzięki czemu kanały i natężenia można porównywać w tym
samym układzie. Po trzecie, skalarne miary udziałów i sektory znaków
zachowują to, co w reprezentacji simpleksowej jest fizycznie
najważniejsze.

Metoda ma również granice. Osie \(\xi_1,\xi_2\) wykorzystują etykiety
trajektorii; przy porównaniu z eksperymentem bez takiej informacji należy
używać miar symetrycznych na permutacje. Dla większej liczby elektronów
rośnie wymiar przestrzeni względnej, więc sam promień \(K\) nie może
zastąpić pełnego rozkładu. Przy porównywaniu różnych natężeń warto
zestawiać \(K\) z bezwymiarowymi miarami, na przykład z \(K/A\), gdzie
\(A=|p_2|+|p_3|+|p_4|\). Wreszcie, przed interpretacją słabych różnic
statystycznych należy dodać przedziały ufności, najlepiej wyznaczone
bootstrapem.

PCA pokazała strukturę jednego trybu kolektywnego i dwóch trybów
względnych. Stała baza Jacobiego--Helmerta przekłada tę obserwację na
układ o jasnej interpretacji fizycznej. Dla \(1{,}7\,\mathrm{PW/cm^2}\)
oba kanały różnią się konsekwentnie: \textit{direct} jest bardziej
kolektywny i równomierny, a \textit{delayed} bardziej zróżnicowany
wewnętrznie. To jest główny wynik pracy.
\end{document}

